\section{Creating a dataset}

Here is a typical sequence of CDI calls used to create a new dataset:

\begin{verbatim}
    gridNew              ! define a grid: from type and size
       ...
    zaxisNew             ! define a Z-axis: from type and size
       ...
    taxisNew             ! define a T-axis: from type
       ...
    vlistNew             ! create a variable list
          ...
       vlistDefVar       ! define variables: from grid and Z-axis
          ...
    streamOpenWrite      ! create a dataset: from name and file type
       ...
    streamDefVlist       ! define variable list
       ...
    streamDefTimestep    ! define time step
          ...   
       streamWriteVar    ! write variable
          ...
    streamClose          ! close the dataset
\end{verbatim}


\section{Reading a dataset}

\begin{verbatim}
    streamOpenRead       ! open existing dataset
       ...
    streamInqVlist       ! find out what is in it
          ...
       vlistInqVarGrid   ! get grid
          ...
       vlistInqVarZaxis  ! get Z-axis
          ...
       vlistInqTaxis     ! get T-axis
          ...
    streamInqTimestep    ! get time step
          ...
       streamReadVar     ! read varible
          ...
    streamClose          ! close the dataset
\end{verbatim}


\section{Compiling and Linking with the CDI library}

Details of how to compile and link a program that uses the CDI C or FORTRAN
interfaces differ, depending on the operating system, the available compilers,
and where the CDI library and include files are installed.
Here are examples of how to compile and link a program that uses the CDI library
on a Unix platform, so that you can adjust these examples to fit your installation.

